\scanpage{483}
\mspage{245}
34. Heuristic profinite description of the category
of algebraic curves \unsure{defined} over \deleted{\unsure{$\mathbb Q$}} finite sub-
extensions $K$ of $\overline{\mathbb Q}_0/\mathbb Q$ (i.e. of $\mathbb C/\mathbb Q$).

We restrict ourselves to \inserted{geometrically} connected curves \inserted{(for convenience)}, anabelian
(by provisional necessity \unsure{!}, cf. below on \unsure{the way} of
getting rid of this restriction). The datum of a
universal covering of \deleted{a} $U_{\overline{\mathbb Q}_0}$ defines an
extension structure
\[
(1)\qquad 1\longrightarrow\pi\longrightarrow\mathcal E\longrightarrow\Gamma^{+}\longrightarrow1\qquad\text{or}
\]
(2) or again we have a homomorphism $\mathcal E\longrightarrow\boldsymbol{\Gamma}$
(image $\Gamma^{+}$ open in $\boldsymbol{\Gamma}$,
kernel called $\pi$),
where $\Gamma^{+}\subset\boldsymbol{\Gamma}_{\mathbb Q}\simeq\boldsymbol{\Gamma}_{0,3}$ is the open \unsure{subgroup|subgroup}
corresponding to $K$. To give such an
extension (provided $\unsure{\mathrm{Centre}}(\pi)=1$) amounts to
giving an outer action of $\Gamma'$ on $\pi$.
A first question: should one put the
loop structure of $\pi$ into the \deleted{\illegible{1}} data \unsure{$\mathcal E$} of
the object (1) (or \unsure{(2)|(1')}) supposed to describe $U\unsure{\mathcal E}/K$?

\uline{Conjecture} It's not necessary -- the loop structure
of $\pi$ is the only loop structure \unsure{invariant}
under the outer action of $\Gamma^{+}$ (or of $\Gamma^{\natural}$). The
loop subgroups $L_i$ are the maximal
subgroups\inserted{in $\pi$}, such that $\mathrm{Norm}_{\mathcal E}(L_i)\longrightarrow\Gamma^{+}$ has
as image an open subgroup of $\unsure{\Gamma^{+}|\Gamma}$.
\margin{\unsure{and} every subgroup $\ne1$ of $\pi$ which has this property is contained in a unique $L_i\ldots$}
I
presume also that every homomorphism \deleted{\unsure{of}}\inserted{between} extensions
$\mathcal E$ of $\Gamma^{+}$ by a $\pi$, $\mathcal E'$ of $\Gamma$ by a $\pi'$
($\mathcal E,\mathcal E'$ coming from algebraic curves $U,U'$) and such
that the image of $\mathcal E$ in $\mathcal E'$ is open, respects
\margin{482}

\scanpage{484}
necessarily the loop structure, and in fact
comes from a homomorphism (unique, as we know) of algebraic
curves.

In any case, if we admit the description of the
loop subgroups, it \unsure{will be} clear that the image by
$\unsure{u}$ of an $L_i$ will be, either $(1)$, or contained
in a unique $L'_j\ldots$

This would mean that the functor from algebraic
curves over $K$\inserted{(morphisms the \unsure{dominant} morphisms)}, to the profinite outer groups
on which $\Gamma$ operates, would be fully faithful.
The functor ``extension of the base field'' from
$K$ to $K'$ corresponds to the functor: restriction
of an \unsure{outer|outer} group (or of an extension
structure) from $\Gamma$ to $\Gamma'$ (open subgroup).
The finite \inserted{\unsure{étale}} coverings of $U$ correspond to the
$\mathcal E$-sets: [$\widetilde U$ \unsure{being chosen}]\ldots

To describe the essential image of this
functor, we are not reduced to conjectures:
we start from the \inserted{canonical} extension $\mathcal E_{03}$ of $\boldsymbol{\Gamma}_{03}=\boldsymbol{\Gamma}$
by $\widehat{\pi}_{03}$, we take an \inserted{open} \unsure{subgroup|subgroup} $\mathcal E'$
of $\mathcal E_{03}$, with image $\Gamma\subset\boldsymbol{\Gamma}_{03}$, kernel
$\pi'\subset\widehat{\pi}_{0,3}$, and in the canonical loop
structure of $\pi'$\inserted{of genus $g$} (induced by that of $\widehat{\pi}_{03}$)
we take an $\unsure{I|\overline I}\subset I(\pi')$\inserted{such that $2g+\unsure{\mathrm{Card}(I)}\geq3$}, \inserted{stable under $\Gamma$, and} we ``plug the holes''
at $I(\pi')\setminus I$. Likewise, to describe when
\margin{483}

\scanpage{485}
\mspage{\unsure{246|245}}
an effective action, lifting an admissible outer
action, is itself admissible -- \deleted{\unsure{can}}
i.e. can be obtained from an algebraic \deleted{\illegible{1}}\inserted{curve}
$U$, \uline{pointed} by a point rational over $K$.
This does not \uline{mean} for all that, that if
$U$ is given by an outer action of $\Gamma$ on a $\pi$,
that the conjugacy classes of admissible liftings
of this action do \unsure{come} all of them from
points of $U$ rational/$K$. But one sees at once that
this will be the case, as soon as one admits
full faithfulness for the \uline{isomorphisms}

\uline{N.B.} Even for the automorphisms of $U_{0,3}$
$=\mathbb P^1_{\mathbb Q}\setminus\{0,1,\infty\}$, this full faithfulness is not
at all clear. One would have to prove that \deleted{\illegible{1}}
the commutant of $\boldsymbol{\Gamma}_{0,3}$ in $\widehat{\widehat{\mathcal T}}_{0,3}$ is reduced
to $\unsure{\mathfrak S_3|\mathcal S_3}$. The situation is less hopeless,
when one gives oneself, with the structure of
\inserted{profinite} outer group with operators $\Gamma$, an arithmetization
of $\pi$ (which supposes that one has \unsure{made explicit}
a loop structure) invariant under \deleted{\illegible{1}} $\Gamma$
-- \unsure{hence} in $\widehat{\widehat{\mathcal T}}(\pi)$ we have at our disposal an $\unsure{N(\pi)|\mathcal N(\pi)}$,
and $\Gamma\hookrightarrow\unsure{N(\pi)|\mathcal N(\pi)}$. In this point of view, \inserted{for} a
homomorphism \inserted{$\pi'\longrightarrow\pi$} of (arithmetized) outer $\Gamma$-groups,
\unsure{it} is understood that it is compatible with
the arithmetizations, i.e. if $\pi''$ is the image of $\pi'$ in $\pi$,
\margin{484}

\scanpage{486}

one sees that the arithmetization of $\pi''$ \unsure{deduced} from that of $\pi'$ by ``passage to the quotient'', be the one \unsure{induced} by $\pi$. \deleted{\unsure{The interest}} In particular, \inserted{for} the automorphisms\inserted{(outer)} of $\pi$ it is understood that \deleted{\illegible{1}} not only do they commute with $\Gamma$, but also that they are in $N(\pi)$. \inserted{(which implies that they are in \unsure{$\widehat{\mathcal T}(\pi)$}, since the centralizer in $\boldsymbol{\Gamma}_\pi$ of \unsure{every} open subgroup of $\boldsymbol{\Gamma}_\pi$ is $(1)$!)} In the case of $U_{0,3}$, we have $N_{0,3}\simeq\unsure{\mathfrak S_3}\times\boldsymbol{\Gamma}_{0,3}$, and we know that the center of $\boldsymbol{\Gamma}_{03}\simeq\boldsymbol{\Gamma}$ is trivial -- so we do find that the group of automorphisms of this structure is reduced \unsure{to|:} \unsure{$\mathfrak S_3$}!

This point of view is nonetheless probably superfluous -- for one \unsure{presumes} that for the given outer action of $\Gamma$, there is a unique arithmetization invariant under $\Gamma$, and that the \inserted{``admissible''} \unsure{homomorphisms} of ``admissible'' outer $\Gamma$-groups, such as they were defined previously, \unsure{automatically} respect this arithmetization. If this were not so, one would of course have to introduce the arithmetizations into the structure.

\deleted{\unsure{Here}, \deleted{\unsure{as} $\unsure{\boldsymbol{\Gamma}_\pi|\Gamma_\pi}\unsure{=}$} \inserted{if $\Gamma=\unsure{\boldsymbol{\Gamma}_{\mathbb Q}}$,} the arithmetization of \deleted{\illegible{1}} $\pi$ \deleted{\unsure{is}} \inserted{must be} \unsure{described} in fact by a well-determined oriented prediscretification $\alpha$, \deleted{\unsure{from the fact that}\illegible{2}} \inserted{In fact, \unsure{we are going to|we ourselves} describe a \uline{strict} \unsure{oriented} prediscretification} \deleted{\illegible{3} \unsure{base point}} well determined of $\pi$, Let us start from an admissible lifting of the given outer operation}

\margin{485}

\scanpage{487}
\mspage{247}

The question arises of finding a more \deleted{\unsure{simple and}} \inserted{simplistic} description of the outer actions of a $\Gamma$ on a $\pi$ which are ``admissible''. \deleted{\illegible{1}} Here, we are going to start from a $\pi$ for which we already \unsure{fix} a loop structure \inserted{and an arithmetization \unsure{a}} (\unsure{left invariant|left invariant} by $\Gamma$), so \deleted{\illegible{2}} that $\Gamma\subset N(\pi)$, $\Gamma\longrightarrow\boldsymbol{\Gamma}_\pi$ \unsure{injective} with open image, \deleted{\illegible{1}} i.e. $\Gamma$ corresponds to a partial splitting (or germ of splitting) of \deleted{\illegible{1}}
\[
1\longrightarrow\widehat{\mathcal T}(\pi)\longrightarrow N(\pi)\longrightarrow\boldsymbol{\Gamma}_\pi\longrightarrow1
\]

\unsure{Let} $(g,\nu)$ be the type of $\pi$; \deleted{\illegible{2}} if $g\geq1$ the simplest criterion, is that the ``\inserted{$\Gamma^{\natural}$-}points'' \unsure{of|of the} $\pi'$ deduced from $\pi$ by plugging all the holes, be distinct. \deleted{\unsure{Except}} (N.B. if $g=1$, \deleted{\unsure{this}} \deleted{\unsure{plugging}} \inserted{by plugging} all the holes we fall into an abelian case -- but that doesn't matter). If $g=0$, \deleted{the criterion is that} there is no condition for $\unsure{\alpha|\nu}=3$, and if $\nu>3$, \unsure{let us set|setting} apart a \deleted{\unsure{a}} subset $I'\subset I(\pi)$ with $\mathrm{card}\,I'=3$, the condition is that on plugging the holes \unsure{of|in} $\unsure{I\setminus I'}$, the \inserted{$\Gamma^{\natural}$-}points \unsure{of the|of} $\pi$ deduced from the points of $\unsure{I\setminus I'}$ are distinct.

Note that in this conjectural perspective, in the limiting cases $(g,0)$ $(g\geq2)$, $(1,1)$, $(0,3)$, we impose no \inserted{a priori} condition on the \unsure{liftings}.

\margin{486}